\chapter*{Abstract}
\label{abstract}

Capacity planning for containerised microservices requires more than measuring CPU utilisation and response time. Classical M/G/$c$ queueing models assume that service times are independent of load, but real servers are executed by operating systems and language runtimes whose scheduling behaviour can change measured wall-clock service time as utilisation grows. This thesis builds a reproducible Docker-based capacity-modelling laboratory and uses it to test when standard queueing assumptions hold for real request-response services.

The experimental platform contains multiple Dockerised server implementations across Apache/PHP, Node.js, Python/Gunicorn, Java, Go, and SQLite-backed Go services. Open-loop Poisson load generators produce request-level traces containing arrival time, service time, queue time, response time, and status code. These traces are replayed through M/G/1 and M/G/$c$ discrete-event simulation (DES) models and validated against measured response-time CDFs using Kolmogorov--Smirnov distance.

The central finding is that several CPU-bound servers exhibit load-dependent wall-clock service time. Mean service time grows with nominal utilisation, violating the constant-service-time assumption used by standard M/G/$c$ models. To model this, the thesis compares M0 (constant service time), M1 (linear growth), and M2 (hyperbolic growth) service-time hypotheses, with M2 given by $E[S](\rho_0)=S_0/(1-\beta\rho_0)$. The model is interpreted as a reduced-form correction for scheduler-induced contention rather than a literal nested Processor-Sharing queue. A further M3 contention-penalty formulation and exploratory priority/feedback DES models are also tested to stress the interpretation.

The evidence chain combines service-time fits, response-CDF simulation, service-time histograms, and ANOVA on $\log(\texttt{service\_ms})$. M2 does not dominate every finite-range case, but it clearly improves over M0 in key workloads where constant service time fails. Structural boundary cases are also identified: Node.js cluster behaves closer to split M/G/1 queues than a shared M/G/3 queue; Go SQLite suggests a tandem application/database-lock queue; and Java DSP is affected by JIT warmup rather than queueing load. The resulting contribution is both pedagogical and research-oriented: a teaching artefact for observing queueing behaviour in real containers, and a mechanistic framework for diagnosing service-time load dependence as a missing factor in container capacity planning.
